% 【本文件写什么】impl 实现层：dummy
%   - 定位：网络设备占位实现。
%   - crate 路径：os/components/wateros-driver/driver-network/network-impl/impl-dummy/src/
%   - 如何实现对应 api-v0 trait：关键类型、状态机、数据结构
%   - 算法或硬件交互流程（可用伪代码/流程图）
%   - 本 impl 启用的 Cargo [features] 与 arch feature（impl-riscv64 等）
%   - 与其它 impl 变体的差异、切换 feature 时的影响
%   - 性能/内存假设与已知限制
% 【事实来源】
%   - 上述 crate 源码与 Cargo.toml
%   - os/feature-tree.txt 中指向本 impl 的 feature 链

\subsubsection{impl-dummy（network）}

\code{DummyNetworkDevice}提供可注册但链路恒down的占位网卡：\code{is\_link\_up}为\code{false}，\code{send}/\code{receive}返回\code{Unsupported}。\code{new\_shared(mac)}直接返回\code{SharedNetworkDevice}便于测试注册表逻辑。

生产bring-up路径不默认注册该设备；\code{stack::init}在无真实网卡时走\code{SmoltcpAdapter::loopback\_only}而非\code{DummyNetworkDevice}。本crate主要用于缺省feature组合下的编译占位与单元测试构造指定MAC的句柄。

\begin{rustcode}
impl NetworkDevice for DummyNetworkDevice {
    fn is_link_up(&self) -> bool { false }
    fn send(&mut self, _buf: &[u8]) -> DriverResult<()> {
        Err(DriverError::Unsupported)
    }
    fn receive(&mut self, _buf: &mut [u8]) -> DriverResult<usize> {
        Err(DriverError::Unsupported)
    }
}
\end{rustcode}
